\section{Proof of Theorem \ref{lemma:approx}}\label{appendix:approx}
In this section, we provide the omitted proof in Section \ref{sec:proof-sketch-approx}.


\subsection{Lemma \ref{lem:covering_division}}
\begin{lemma}[Lemma 20 in \citet{nakada2020adaptive}]\label{lem:covering_division}
    Let $\cC=\{I_k\}_{k=1}^{N_r(S)}$ be a minimum $r$-covering of $S$ where $I_k$'s are hypercubes with side length $r$.  Then, there exists a disjoint partition $\{\cC_j\}_{j=1}^{5^d} \subset \cC$ such that $\cC = \bigcup_{j=1}^{5^d} \cC_j$ and $d(I_i, I_l) \geq r$ hold for any $I_i \neq I_l \in \cC_j$ if $\card(\cC_j) \geq 2$, where $d(\cA,\cB):= \inf \{\norm{x-y} | x\in\cA,y\in\cB\}$ is defined as the distance of any two sets $\cA$ and 
    $\cB$.  % for all $i = 1, \dots, 5^D$.
\end{lemma}

\subsection{Realization of hypercube determination function $g_I^{\text{ind},r}$}\label{sec:realize-ind}

Hypercube determination function $g_I^{\text{ind},r}$ can be realized by weight matrices and intercepts $(4,0)\odot(W^2,-d)\odot[(W_1^1,b_1^1),\ldots,(W_d^1,b_d^1)]$ where $W_i^1,b_i^1$ and $W^2$ are defined by
\begin{align*}
    W_i^1 := \begin{pmatrix}e_i^\top & e_i^\top & e_i^\top & e_i^\top \\ 0 & 0 & 0 & 0\end{pmatrix}^\top,~~
    b_i^1 := \biggl(-\iota_i + r ~~ -\iota_i + \frac{r}{2} ~~ -\iota_i - \frac{r}{2} ~~ -\iota_i - r \biggr),
\end{align*}
and
\begin{align*}
    W^2 = \bigg(&\underbrace{\frac{2}{r}, -\frac{2}{r}, -\frac{2}{r}, \frac{2}{r},  \frac{2}{r}, -\frac{2}{r}, -\frac{2}{r}, \frac{2}{r}, \dots, \frac{2}{r}, -\frac{2}{r}, -\frac{2}{r}, \frac{2}{r}}_{4d},\frac{1}{4} \bigg).
\end{align*}
The above realization gives exactly the form in \eqref{eq:identical}. Moreover, we summarize the properties of $g_I^{\text{ind},r}$ as following:
\begin{proposition}\label{prop:ind}
    For any $x\in \cD_{\rm data}$ and $y \in \RR$, we have
    \begin{align*}
        g_I^{\text{ind},r}(x,y)
        \begin{cases}
            =y,  &x\in I \text{ and } y \in [0,4], \\
            \leq y,  &x\in I\oplus \frac{r}{2}\text{ and } y \in [0,4], \\
            =0,  &\text{otherwise}. \\
        \end{cases}
    \end{align*}
    Furthermore, we obtain the following properties
    \begin{enumerate}
        \item $L(g_I^{\text{ind},r})=3$,
        \item $B(g_I^{\text{ind},r})\leq \max\{4,d,1+r,2/r\}$,
        \item $K(g_I^{\text{ind},r})=24d+6$.
    \end{enumerate}
\end{proposition}

%\subsection{Proof of Lemma \ref{lemma:poly}}
%\begin{proof}
%Firstly, we rewrite the target polynomial $\sum_{\abs{\alpha} < \beta} c_{k, \alpha} (x - x_k)^\alpha$.
%By the binomial theorem \citet{folland2013real}, we have
%\begin{equation*}
%    (x - x_k)^\alpha = \sum_{\gamma \leq \alpha} \binom{\alpha}{\gamma} (-x_k)^{\alpha - \gamma} x^\gamma.
%\end{equation*}
%Then, we evaluate the polynomial as
%\begin{align*}
%    \sum_{|\alpha| < \beta} c_{k, \alpha} (x - x_k)^\alpha &= \sum_{|\alpha| \leq \lfloor \beta \rfloor} \Biggl\{ \sum_{\gamma \leq \alpha} c_{k, \alpha} \binom{\alpha}{\gamma} (-x_k)^{\alpha - \gamma} x^\gamma \Biggr\}\\
%    &= \sum_{|\gamma| \leq \lfloor \beta\rfloor} \Biggl\{ \sum_{\gamma \leq \alpha, |\alpha| \leq \lfloor \beta \rfloor} c_{k, \alpha} \binom{\alpha}{\gamma} (-x_k)^{\alpha - \gamma} \Biggr\} x^\gamma\\
%    &=: \sum_{|\gamma| \leq \lfloor \beta\rfloor} \tilde c_{k, \gamma} x^\gamma.
%\end{align*}
%Since $\binom{\alpha}{\gamma} \leq \alpha_1^{\gamma_1} \dots \alpha_d^{\gamma_d}$, we can bound $\abs{\tilde c_{k, \gamma}}$ by
%\begin{align*}
%    |\tilde c_{k, \gamma}| &\leq \sup_{x \in S} \Biggl|\sum_{|\alpha| \leq \lfloor \beta \rfloor} \sum_{\gamma \leq \alpha} c_{k, \alpha} \binom{\alpha}{\gamma} (-x_k)^{\alpha-\gamma} \Biggr|\\
%    &\leq  \sum_{|\alpha| \leq \lfloor \beta \rfloor} \sum_{\gamma \leq \alpha} \alpha_1^{\gamma_1} \dots \alpha_d^{\gamma_d} R_S^\beta \\
%    &\leq  \sum_{|\alpha| \leq \lfloor \beta \rfloor} (\alpha_1+1)\dots(\alpha_d+1) \alpha_1^{\alpha_1} \dots \alpha_d^{\alpha_d} R_S^\beta\\
%    &\leq  d^{\lfloor \beta \rfloor+1} (1+\lfloor \beta \rfloor)^{d(1 +\lfloor \beta \rfloor)} R_S^\beta.
%\end{align*}
%\end{proof}





\subsection{Bounding the approximation error}\label{sec:approx-error-bound}

Firstly, we compute the approximation error of using $g^{f_0}_\epsilon$ to estimate $f_0=f^*+2$. Recall that we defined $g^{f_0}_\epsilon:=g^{\max,5^d} \circ g^{\text{sum}} \circ g^{\text{simul}}_{\epsilon/2}$. When $x \in S$, there exists some $I \in \cC$ such that $x \in I$. Then for this $x$, we have
\begin{align*}
    g^{f_0}_\epsilon(x) &= \max_{j \in [5^d]} \sum_{I_k \in \cC_j} g^{\text{simul}}_{\epsilon/2,k}(x)
    = \max_{I_k \in \neig(I)}  g^{\text{simul}}_{\epsilon/2,k}(x)  
    \leq \max_{I_k \in \neig(I)}  g^{\poly}_{\epsilon/2,k}(x),
\end{align*}
where $\neig(I)=\{ I' \in \cC | (I \oplus 3r/2)\cap I' \neq \emptyset \}$ denotes the $3r/2$-neighborhood of hypercube $I$. 
In other words, when computing $g^{f_0}_\epsilon(x)$, we only need to take maximum over the estimated function value within hypercubes near $x$. The second equality follows from the Proposition \ref{prop:ind} that 
 $g^{\text{simul}}_{\epsilon/2,l}(x)=0$ for $I_l \not \in \neig(I)$ and $d(I_l, I_k) > r$ holds for $I_l \neq I_k \in \cC_i$ for all $i$. The last inequality is due to the construction of $g^{\text{simul}}_{\epsilon/2}$ in \eqref{eq:simul}.

Given $\epsilon \in (0, 1)$, we ensure $0 \leq g^{\text{simul}}_{\epsilon/2,k}(x) \leq 4$ for all $I_k \in \cC$ by Proposition \ref{prop:ind}, since $g^{\text{simul}}_{\epsilon/2,k}$ approximates $f_{I_k}$ which is an $\epsilon/2$-accuracy Taylor polynomial of $f_0 \in [1, 3]$. When $x \in I_t$, the error is bounded as
\begin{align*}
    | g^{f_0}_\epsilon (x) - f_0(x) | 
    =& \max \Bigl\{  \max_{I_k \in \neig(I_t)}  g^{\text{simul}}_{\epsilon/2,k}(x) -f_0(x), f_0(x)- \max_{I_k \in \neig(I_t)}  g^{\text{simul}}_{\epsilon/2,k}(x)         \Bigr\} \\
    \leq & \max \Bigl\{  \max_{I_k \in \neig(I_t)}  g^{\poly}_{\epsilon/2,k}(x)-f_0(x), f_0(x)- g^{\text{poly}}_{\epsilon/2,t}(x) \Bigr\}\\
    \leq& \max_{I_k \in \neig(I_t)}  \Bigl| g^{\poly}_{\epsilon/2,k}(x)-f_0(x) \Bigr| \\
    \leq& \max_{I_k \in \neig(I_t)}  \Bigl| g^{\poly}_{\epsilon/2,k}(x)-f_{I_k}(x) \Bigr|
     + \max_{I_k \in \neig(I_t)}  \bigl| f_{I_k}(x) -f_0(x) \bigr| \\
    \leq& \frac{\varepsilon}{2} + d^\beta \biggl(\frac{3r}{2} \biggr)^\beta \leq \epsilon,
\end{align*}
where the last inequality follows from \eqref{eq:poly} and Lemma \ref{lemma: taylor_approx}. 
In terms of parameter tuning, we choose $r = d^{-1} \epsilon^{1/\beta}/2$. 

Next, we extend approximation results of $f_0$ to $f^*$. To do so, we firstly implement a neural network  $g^{\text{mod}}(z)=(-z+1)\circ\relu(-z+2)\circ \relu(z-1)$, which has the equivalent form $g^{\text{mod}}(z)=\min(\max(1,x),3)-2$ for any $z\in \RR$. In addition, $g^{\text{mod}}$ has the following properties:
\begin{equation*}
    L(g^{\text{mod}}) = 3,~~B(g^{\text{mod}}) \leq 2,~~\text{and  }K(g^{\text{mod}}) = 12.
\end{equation*}

Then consider $g^{f^*}_\epsilon = g^{\text{mod}}\circ g^{f_0}_\epsilon$ to obtain the desired approximation error $\epsilon$ for any $x\in S$ 
\begin{align*}
    \sup_{x \in \cD_{\rm data}} |g^{f^*}_\epsilon(x) - f^*(x)| 
    =& \sup_{x \in \cD_{\rm data}} \bigl|\min\bigl(\max\bigl(1,g^{f_0}_\epsilon(x) \bigr),3\bigr) - (f^*(x) + 2)\bigr|\\
    =& \sup_{x \in \cD_{\rm data}} \bigl|\min\bigl(\max\bigl(1,g^{f_0}_\epsilon(x) \bigr),3\bigr) - f_0(x) \bigr|\\
    \leq& \sup_{x \in \cD_{\rm data}} \bigl| g^{f_0}_\epsilon(x)  - f_0(x) \bigr|\\
    \leq& \epsilon.
\end{align*}





\subsection{Computing network sizes}\label{sec:parameter}

Recall that the ReLU neural network $g^{f^*}_\epsilon$ is defined as  
\begin{align*}
    g^{f^*}_\epsilon=& g^{\text{mod}}\circ g^{\max,5^d} \circ g^{\text{sum}} \circ g^{\text{simul}}_{\epsilon/2}\\
    =& g^{\text{mod}}\circ g^{\max,5^d} \circ g^{\text{sum}} \circ\bigl(g_{I_1}^{\text{ind},r}\circ g^{\text{filter}}_1 ,\ldots,g_{I_{ N_r(S)}}^{\text{ind},r}
    \circ g^{\text{filter}}_{ N_r(S)} \bigr) 
    \circ (g^{Id}_{d,L},  g^{\poly}_{\epsilon/2}).
\end{align*}
Note that $N_r(S) \leq c_0 r^{-p}$.
Combined with sub-neural network structures given in Appendix B.1.1 of \citet{nakada2020adaptive}, $g^{f^*}_\epsilon$ has the following properties:
\begin{align*}
    L(g^{f^*}_\epsilon(x)) &= L(g^{\text{mod}}) + L(g^{\max,5^d}) + L(g_{I_1}^{\text{ind},r}) + L(g^{\text{filter}}_1) + L(g^{\poly}_{\epsilon/2})\\
    &\leq 11 + 2d \log_2 5 + (11 + (1+\beta)/d)(2 + \log_2\beta),\\
    B(g^{f^*}_\epsilon(x)) &\leq \max\bigl\{B(g^{\text{mod}}), B(g^{\max,5^d}),B(g_{I_1}^{\text{ind},r}),B(g^{\text{filter}}_1) , B(g^{Id}_{d,L}),B(g^{\poly}_{\epsilon/2})\bigr\},\\
    &\leq \max\{d, 1 + r, 2 /r, c_1^{\poly} R_S^{\beta s_1^{\poly}}\epsilon^{-s_1^{\poly}} \},\\
    &\leq \max\{ 4d \epsilon^{-1/\beta} , c_1^{\poly} R_S^{\beta s_1^{\poly}}\epsilon^{-s_1^{\poly}} \},\\
    K(g^{f^*}_\epsilon(x)) &\leq 2K(g^{\text{mod}}) + 2K(g^{\max,5^d})  + 2N_r(S)\cdot K(g_{I_1}^{\text{ind},r}\circ g^{\text{filter}}_1) + 2 K(g^{Id}_{d,L}) + 2K(g^{\poly}_{\epsilon/2})\\
    &\leq  2 c_2^{\poly} R_S^p \epsilon^{-p/\beta} + 2(50d + 17 + c_2^{\poly})   N_r(S)\\
    &\quad+ 2(12 + 42 \times 5^d + 2d + 2d(11 + (1+\beta)/p)(2 + \log_2 \beta))),\\
    &\leq  2 c_2^{\poly} R_S^p \epsilon^{-p/\beta} + 2(50d + 17 + c_2^{\poly})c_0 (2d)^p\epsilon^{-p/\beta}  \\
    &\quad+ 2(12 + 42 \times 5^d + 2d + 2d(11 + (1+\beta)/p)(2 + \log_2 \beta))),
\end{align*}
where $c_2^{\poly}=O(d^{2+\lfloor\beta\rfloor})$. By adjusting several constants, we obtain the statement.




\section{Proof of Theorem \ref{thm:generalization}}\label{sec:proof-generalization}

The proof of Theorem \ref{thm:generalization} mainly follows \citet{chen2022nonparametric}. Our Lemma \ref{lemma:t1} and \ref{lemma:t2} is a slight revision of Lemma 5 and 6 in \citet{chen2022nonparametric}, where we substitute $\ell_\infty$ covering number with $\ell_2$ covering number to deal with unbounded domain. In this section, we compute the $\ell_2$ covering number of neural network class and then present the proof of  Theorem \ref{thm:generalization}.

\subsection{Proof of Lemma \ref{lemma:covering-number}}
\begin{proof}[Proof of Lemma \ref{lemma:covering-number}]
    To construct a covering for $\cF(H,L,B,K)$, we discretize each parameter by a unit grid with grid size $h$. Recall that we write $f \in \cF(H,L,B,K)$ as $f(x) = W_L \cdot \relu (W_{L-1}\cdots \relu(W_1 x +b_1) \cdots +b_{L-1}) +b_L$ in \eqref{eq: nn-form}. Choose any $f,f' \in \cF(H,L,B,K)$ with parameters at most $h$ apart from each other. Denote the weight matrices and intercepts in $f,f'$ as $W_L,\ldots,W_1,b_L,\ldots,L_1$ and $W_L',\ldots,W_1',b_L',\ldots,L_1'$ respectively, where $W_l \in \RR^{{d_l}\times d_{l-1}}$ and $b_l \in \RR^{d_l}$ for $l=1,\ldots,L$. Without loss of generality, we assume $d_l \leq K$ since all the parameters have at most $K$ nonzero entries. If the input dimension is larger than $K$, we let the redundant dimensions of input equal to zeros.
    
    Notice that for any random variable $y \in \RR^{d_{l-1}}$ which is subject to a distribution $P_Y$, we have
    \begin{equation*}
         \int_{\RR^{d_{l-1}}} \bignorm{ ( W_l y+ b_l)-( W_l' y+ b_l')}_2^2 \ud P_Y(y) =  \int_{\RR^{d_{l-1}}}   \bignorm{ \sum_{i=1}^{d_{l-1}} (W_{l,i} - W_{l,i}') y_i+(b_l- b_l') }_2^2 \ud P_Y(y).
    \end{equation*}
    By the inequality $\norm{t+s}^2 \leq 2\norm{t}^2 + 2\norm{s}^2$ which holds for any $s,t\in \RR^{d_l}$, we obtain
    \begin{align*}
        \int_{\RR^{d_{l-1}}} \bignorm{ ( W_l y+ b_l)-( W_l' y+ b_l')}_2^2 \ud P_Y(y) 
        \leq & 2 \int_{\RR^{d_{l-1}}}  \bignorm{\sum_{i=1}^{d_{l-1}}  (W_{l,i}-W_{l,i}') y_i  }^2 \ud P_Y(y)  + 2 \norm {b_l-b_l'}_2^2 \\
        \leq &  2 \sup_{i=1\ldots,d_{l-1}} \norm{ W_{l,i}-W_{l,i}' }_2^2 \cdot \int_{\RR^{d_{l-1}}}  \sum_{i=1}^{d_{l-1}} y_i^2  \ud P_Y(y)  + 2 \norm {b_l-b_l'}_2^2.
    \end{align*}
    Since parameters $W_l,b_l$ differ at most $h$ from $W_l',b_l'$ with respect to each entry, we get 
    \begin{align}\label{eq:difference1}
    \int_{\RR^{d_{l-1}}} \bignorm{ ( W_l y+ b_l)-( W_l' y+ b_l')}_2^2 \ud P_Y(y) 
        &\leq  2 d_{l-1} h^2 \norm{y}^2_{L_2(P_Y)} + 2 d_l h^2 \notag \\
        &\leq   2 K h^2 \norm{y}^2_{L_2(P_Y)} + 2 K h^2 .
    \end{align}
Similarly, we have
\begin{align}\label{eq:difference2}
    \int_{\RR^{d_{l-1}}} \bignorm{  W_l y+ b_l }_2^2 \ud P_Y(y) = & \int_{\RR^{d_{l-1}}}   \bignorm{ \sum_{i=1}^{d_{l-1}} W_{l,i}  y_i + b_l }_2^2 \ud P_Y(y) \notag \\
        \leq & 2 \int_{\RR^{d_{l-1}}}  \bignorm{ \sum_{i=1}^{d_{l-1}}  W_{l,i} y_i}^2 \ud P_Y(y)+ 2 \norm {b_l}_2^2 \notag  \\
        \leq &  2 \sup_{i=1\ldots,d_{l-1}} \norm{ W_{l,i} }_2^2 \cdot \int_{\RR^{d_{l-1}}}  \sum_{i=1}^{d_{l-1}} y_i^2  \ud P_Y(y)  + 2 \norm {b_l}_2^2 \notag \\
        \leq &  2 d_{l-1} B^2 \norm{y}^2_{L_2(P_Y)}  + 2 d_l B^2 \notag\\
        \leq &  2 K B^2 \norm{y}^2_{L_2(P_Y)}  + 2 K B^2.
\end{align}

Since the ReLU actiavtion function is $1$-Lipschitz continuous for each coordinate, we can apply \eqref{eq:difference1} and \eqref{eq:difference2} repeatedly to bound $ \norm{f-f'}^2_{L_2(P_{\rm data},S)}$:
\begin{align*}
  \norm{f-f'}^2_{L_2(P_{\rm data},S)}
  =& \int_{S}  \bigl|  W_L \cdot \relu (W_{L-1}\cdots \relu(W_1 x +b_1) \cdots +b_{L-1}) +b_L \\
  &\quad -  W_L' \cdot \relu (W_{L-1}'\cdots \relu(W_1' x +b_1') \cdots +b_{L-1}') -b_L'    \bigr|^2 \ud P_{\rm data}(x) \\
  \leq& 2 \int_{S}   \bigl|  W_L \cdot \relu (W_{L-1}\cdots \relu(W_1 x +b_1) \cdots +b_{L-1}) +b_L \\
  &\quad -  W_L' \cdot \relu (W_{L-1}\cdots \relu(W_1 x +b_1) \cdots +b_{L-1}) -b_L'  \bigr|^2 \ud P_{\rm data}(x)\\
  &\quad + 2 \norm{W_L'}_2^2 \int_{S}   \bigl\|   (W_{L-1}\cdots \relu(W_1 x +b_1) \cdots +b_{L-1})  \\
  &\quad -   (W_{L-1}' \cdots \relu(W_1' x +b_1') \cdots +b_{L-1}')  \bigr\|_2^2 \ud P_{\rm data}(x)\\
  \leq & 4 K h^2+ 4 K h^2  \int_{S}   \norm{  W_{L-1}\cdots \relu(W_1 x +b_1) \cdots +b_{L-1} }^2 \ud P_{\rm data}(x)\\ 
  &\quad +  2 K B^2 \int_{S}   \bigl\|   (W_{L-1}\cdots \relu(W_1 x +b_1) \cdots +b_{L-1})  \\
  &\quad -   (W_{L-1}' \cdots \relu(W_1' x +b_1') \cdots +b_{L-1}')  \bigr\|_2^2 \ud P_{\rm data}(x).
\end{align*}
Besides, we derive the following bound on $\norm{ W_{L-1}\cdots \relu(W_1 x +b_1) \cdots +b_{L-1}}_{L^2(P_{\rm data},S)}$:
\begin{align*}
    \int_{S}   &\norm{  W_{L-1}\cdots \relu(W_1 x +b_1) \cdots +b_{L-1} }^2 \ud P_{\rm data}(x) \\
    &\quad \leq  2KB^2 \int_{S}   \norm{  W_{L-2}\cdots \relu(W_1 x +b_1) \cdots +b_{L-2} }^2 \ud P_{\rm data}(x) +2KB^2 \\
    &\quad \leq (2KB^2)^{L-1} d R_S^2 +  (2KB^2)^{L-1} \\
    &\quad \leq 2^L(KB^2)^{L-1} d R_S^2,
\end{align*}
where the last inequality is derived by induction and $\|x \|^2 =\sum_{i=1}^d  x_i^2 \leq d R_S^2$ for any $x\in S$. Substituting back into the bound for $ \norm{f-f'}^2_{L_2(P_{\rm data},S)}$ , we obtain
\begin{align*}
    \norm{f-f'}^2_{L_2(P_{\rm data},S)} \leq& 4 K h^2+ 2^{L+2} K^L B^{2(L-1)} h^2 d R_S^2 \\
    & \quad+  4 K B^2 \int_{S}   \bigl\|   (W_{L-1}\cdots \relu(W_1 x +b_1) \cdots +b_{L-1})  \\
    &\quad -   (W_{L-1}' \cdots \relu(W_1' x +b_1') \cdots +b_{L-1}')  \bigr\|_2^2 \ud P_{\rm data}(x)\\
    \leq & 4 (L-1) (K B^2)^{L-1} h^2+ 2^{L+2} (L-1) K^L B^{2(L-1)} h^2 d R_S^2 \\
    & \quad + ( 2 K B^2)^{L-1} \int_{S}   \norm{ W_1 x +b_1  -   W_1' x -b_1' }_2^2 \ud P_{\rm data}(x)\\
    \leq & 4^{L-1} L K^L B^{2L} h^2 d R_S^2,
\end{align*}
where the second inequality is obtained by induction. Therefore, combining the above inequality with $\|f\|_\infty \leq 1$ for any $f \in\cF(L,B,K)$ and $P_{\rm data}(S^c) \leq \tau$, we get
\begin{align*}
    \norm{f-f'}^2_{L_2(P_{\rm data})} =& \int_{S}  | f(x) - f'(x) |^2  \ud P_{\rm data}(x) +  \int_{S^c}  | f(x) - f'(x) |^2  \ud P_{\rm data}(x) \\
    \leq &  4^{L-1} L K^L B^{2L-2} h^2 d R_S^2 + 4\tau.
\end{align*}
Now we choose $h$ satisfying $h=\sqrt{(\delta^2-4\tau)/(4^{L-1}L K^L B^{2L-2} d R_S^2)}$. Then discretizing each parameter uniformly into $2B/h$ grid points yields a $\delta$-covering on $\cF(L, B,K)$.  Moreover, the covering number $\cN_2(\delta, \cF(L, B,K))$ satisfies
\begin{align*}
    \log \cN_2(\delta, \cF( L, B,K)) \leq K \log\biggl( \frac{2B}{h} \biggr) = K \log\biggl( \frac{2^L \sqrt{dL} K^{L/2} B^L R_S }{\sqrt{\delta^2-4\tau}} \biggr).
\end{align*}
\end{proof}


\subsection{Proof of Theorem \ref{thm:generalization}}
\begin{proof}[Proof of Theorem \ref{thm:generalization}]
The square error of the estimator $\hat{f}$ can be decomposed into a squared bias term and a variance term, which can be bounded using the covering number of the function class. According to Lemmas \ref{lemma:t1} and \ref{lemma:t2}, for any constant $\delta\in(0,1)$, we have
\begin{equation}\label{eq: l2-bound}
    \begin{split}
        \EE \| \hat{f} - f^*\|^2_{L^2(P_{\rm data})} \leq& c \inf_{f \in \cF( L, B, K)} \norm{ f(x) - f^*(x)}^2_{L^2(P_{\rm data})}  + c \sigma^2 \frac{\log \cN_2(\delta, \cF( L, B, K)) + 2}{n} \\
        &\quad +  c\bigg(   \sqrt{\frac{\log \cN_2(\delta, \cF(L, B, K)) + 2}{n}} + 1 \bigg)\sigma \delta + \frac{c'}{3n} \log \cN_2(\delta/4H, \cF(L, B, K))+ c'\delta.
    \end{split}
\end{equation}

%Next, by Assumption \ref{assumption:pdata_cover}, for any sufficiently small $r,\tau>0$, there exists $p=p(r,\tau)$  and $S\subseteq [0,1]^d$ satisfying $P_{\rm data}(S) \geq 1-\tau$ and $N_r(S) \leq c_0 r^{-p}$. Besides, notice that for any $f \in \mathcal{F}(L,B,K)$, the square error of $f$ can be divided into one part of the high-density region $S_\epsilon$ and the other of the rest area $S^c$. Thereby we get 
%\begin{align*}
 %   \int_{\RR^d} (f(x)-f^*(x))^2 \ud P_{\rm data}(x) =&  \int_{S_\epsilon^c} (f(x)-f^*(x))^2 \ud P_{\rm data}(x) +\int_{S_\epsilon} (f(x)-f^*(x))^2 \ud P_{\rm data}(x) \\
 %   \leq& 4\tau + \int_{S_\epsilon} (f(x)-f^*(x))^2 \ud P_{\rm data}(x) ,
%\end{align*}
%where the last inequality is given by $P_{\rm data}(S^c) \leq \tau$. 
Choose $\epsilon=(2dr)^\beta $ in Theorem \ref{lemma:approx}.
Accordingly, we set tuple $(L,B,K)$ as
\begin{equation*}
     L=C_1,~~ B=O(R_S^{\beta s} r^{-s\beta}),~~ \text{and}~ K=C_2 R_S^p r^{-p}.
\end{equation*}
Then we have
%there exists a network class $\mathcal{F}(L,B,K)$ that yields a function $f$ approximating $f^*$ up to an error $\epsilon$ on $S$, i.e. 
\begin{equation*} 
    \inf_{f \in \cF( L, B, K)} \norm{ f(x) - f^*(x)}^2_{L^2(P_{\rm data})} \leq (2dr)^{2\beta} + 4\tau. 
\end{equation*} 
Invoking the upper bound of the covering number in Lemma \ref{lemma:covering-number}, we derive
\begin{align*}
    \EE \norm{ \hat{f} - f^*}^2_{L^2(P_{\rm data})} \leq& (2dr)^{2\beta} + 4\tau +  \frac{c\sigma^2}{n}  \bigg(  K \log\biggl( \frac{2^{L} \sqrt{dL} K^{L/2} B^L R_S }{\sqrt{\delta^2-4\tau}} \biggr) +2c \bigg) \\
        &\quad + c\sqrt{\frac{ K \log(2^{L} \sqrt{dL} K^{L/2} B^L R_S /\sqrt{\delta^2-4\tau}) +1}{n}} \sigma\delta \\
        &\quad + \frac{c'}{3n}   K \log\biggl( \frac{2^{L} \sqrt{dL} K^{L/2} B^L R_S }{\sqrt{\delta^2-4\tau}}  \biggr) +(c\sigma +c')\delta \\
        =& O\bigg( \tau + d^{2\beta} r^{2\beta} +  \frac{\sigma^2 K}{n} \log\biggl( \frac{ \sqrt{dL} K^{L/2} B^L R_S }{\sqrt{\delta^2-4\tau}}  \biggr) \\
        &\quad + \biggl(  \frac{K \log (\sqrt{dL} K^{L/2} B^L R_S /\sqrt{\delta^2-4\tau} ) }{n} \biggr)^{1/2}  \sigma \delta+ \sigma \delta+ \frac{1}{n}\bigg).
\end{align*}
By Cauchy-Schwartz inequality, for $0<\delta<1$, we have
\begin{align*}
    \EE \norm{ \hat{f} - f^*}^2_{L^2(P_{\rm data})} 
        =& O\biggl( \tau + d^{2\beta} r^{2\beta} +  \frac{\sigma^2 K}{n} \log\biggl( \frac{ \sqrt{dL} K^{L/2} B^L R_S }{\sqrt{\delta^2-4\tau}}  \biggr)+\sigma\delta+ \frac{1}{n}\biggr).
\end{align*}
Plugging in our choice of $(L,B,K)$, we get
\begin{align*}
    \EE \norm{ \hat{f} - f^*}^2_{L^2(P_{\rm data})} 
    = & O\biggl( \tau + d^{2\beta} r^{2\beta} + \frac{\sigma^2 d}{n}\biggl(\frac{R_S}{r} \biggr)^p \log\biggl( \frac{(R_S/r)^p}{ \delta^2-4\tau}\biggr)+ \sigma \delta + \frac{1}{n}\biggr)
\end{align*}
Now we choose  $\delta=r^{2\beta}$. Then we deduce the desired estimation error bound
\begin{align*}
    \EE \| \hat{f} - f^*\|^2_{L^2(P_{\rm data})} &= O\biggl(   \tau +  \sigma r^{2\beta} + \frac{\sigma^2}{n}  \biggl(\frac{R_S}{r} \biggr)^p \log\biggl( \frac{(R_S/r)^p}{ r^{4\beta}-4\tau}\biggr) + \frac{1}{n}\biggr)\\
    & = O\biggl(   \tau +  \sigma r^{2\beta} + \frac{\sigma^2}{n}  \biggl(\frac{R_S}{r} \biggr)^p \log\biggl( \frac{(R_S/r)^p}{ r^{4\beta}-4\tau}\biggr)\biggr).
\end{align*}
The last equality is due to $R_S/r>1$.
%where constant $C$ depends on $R,\beta,d,p,s$ and $\sigma$.
%Denote $f_\epsilon$ as the function that approximates $f^*$ up to an error $\epsilon$ on $S_\epsilon$, i.e, $\sup_{x \in S_\epsilon} |f_\epsilon(x)-f^*(x)| \leq \epsilon$, whose existence is proven in Lemma \ref{lemma:approx}. Then we can decompose the mean square error as
%\begin{align*}
 %   \| \hat{f} - f^*\|^2_{L^2(P_{\rm data})} \leq&  2\| \hat{f} - f_\epsilon\|^2_{L^2(P_{\rm data})} + 2 \| f_\epsilon - f^*\|^2_{L^2(P_{\rm data})}\\
 %   \leq& 2\| \hat{f} - f_\epsilon\|^2_{L^2(P_{\rm data})} + 2\int_{S_\epsilon} (  f_\epsilon (x)-f^*(x))^2 \ud(P_{\rm data}(x)) + 8R^2 \epsilon^2,
%\end{align*}
%where the last inequality follows from $P_{\rm data}(S_\epsilon^c) \leq \epsilon^2$. By Theorem 7 of \cite{nakada2020adaptive}, with probability at least $1-2\exp(-n^{p/(2\beta+p)})$,
%\begin{align*}
%    \| \hat{f} - f_\epsilon\|^2_{L^2(P_{\rm data})} \leq& \bigg( 2^{30}\sigma^4+10+\frac{2^8 M^2}{3} \bigg) n^{-2\beta/(2\beta+p)} + 72 \frac{1}{n} \sum_{i=1}^n (f_\epsilon(x_i)-f^*(x_i))^2 \\
%    & \quad + 2^{17} \sigma^2 \frac{K}{n} \Big( \log \Big( L(K+1)^L B^L n^{\beta/(2\beta+p)}+1 \Big) \Big)^2\\
%    &\quad + \frac{2^8 M^2 K}{3n} \log \Big( 2n^{\beta/(2\beta+d)}L(K+1)^L B^L \Big).
%\end{align*}
%Furthermore, notice that
%\begin{align*}
%    \frac{1}{n} \sum_{i=1}^n (f_\epsilon(x_i)-f^*(x_i))^2 =& \frac{1}{n} \sum_{x_i \in S_\epsilon} (f_\epsilon(x_i)-f^*(x_i))^2 + \frac{1}{n} \sum_{x_i \notin S_\epsilon} (f_\epsilon(x_i)-f^*(x_i))^2 \\
%    \leq & \epsilon^2 +  \frac{n_c}{n} R^2 ,
%\end{align*}
%where $n_c = \sum_{i=1}^n \II(x_i \notin S_\epsilon) $ is the number of samples out of region $S_\epsilon$ and subject to binomial distribution $\mathrm{B}(n,\epsilon^2)$. 
\end{proof}


\section{Proof of Lemma \ref{lemma:gaussian}}\label{proof:gaussian}

\begin{proof}[Proof of Lemma \ref{lemma:gaussian}]
    For simplicity of proof, set $Q=I$. By the construction of $S(R,r;p)$ in \eqref{eq:set}, we notice that
    \begin{align*}
        \PP(X \notin S(R,r;p) ) \leq & \PP\bigg(\sum_{i=1}^p \frac{x_i^2}{\lambda_i^2} >  R^2 \bigg) + \PP\bigg(|x_j|>\frac{r}{2} \text{ for some } j \in [p+1,d]\bigg).
    \end{align*}
    Since $X_{1:p}=(x_1,\ldots,x_p)\sim N(0,\Lambda_p)$ where $\Lambda_p=\diag(\lambda_1^2,\ldots,\lambda_p^2)$, the variable $Z=\Lambda_p^{-1/2} X_{1:p} \sim N(0,I_p)$ is a standard Gaussian. Then for any fixed $R>0$, the probability $\PP(\sum_{i=1}^p x_i^2/\lambda_i^2 >  R^2 )
    $ is equal to $ \PP ( \|Z\|^2 > R^2 )$. Moreover, by Lemma \ref{lemma:gaussian-tail}, if we choose $R^2>p$, we will have
    \begin{align*}
         \PP\big( \|Z\|^2 > R^2 \big) \leq&  \bigg( \frac{2R^2 +p}{p} \bigg)^{\frac{p}{2}} \exp \bigg(-\frac{R^4}{2R^2+p} \bigg) \\
         =&  \exp \bigg(-\frac{p}{2} \cdot \frac{R^4/p^2}{R^2/p+1/2} + \frac{p}{2} \cdot \log\bigg( \frac{2R^2}{p}+1 \bigg) \bigg) \\
         =& O \Bigg( \bigg[ \exp \bigg(-\frac{2R^2}{3p} + \log\bigg( \frac{3R^2}{p} \bigg) \bigg) \bigg]^{\frac{p}{2}} \Bigg).
    \end{align*} 
    Besides, by Lemma \ref{lemma:gaussian-tail-1}, for $j=p+1,\ldots,d$, we derive 
    \begin{align*}
        \PP\bigg(|x_j|>\frac{r}{2} \bigg) = \PP\bigg( \bigg| \frac{x_j}{\lambda_j} \bigg|>\frac{r}{2\lambda_i} \bigg) 
        =  O\bigg( \exp\bigg\{ -\frac{r^2}{8\lambda_j^2} \bigg\} \bigg).
    \end{align*}
    Then we can apply the union bound of probability to get
    \begin{align*}
        \PP\bigg(|x_j|> \frac{r}{2} \text{ for some } j \in [p+1,d]\bigg) \leq \sum_{j=p+1}^d  \PP\bigg(|x_j| > \frac{r}{2} \bigg) 
        =  O\bigg( \sum_{j=p+1}^d \exp\bigg( -\frac{r^2}{8\lambda_j^2} \bigg) \bigg).
    \end{align*}
    Recall that we choose $\lambda_p^{-1}=2R/r$. Then we have
    \begin{align*}
        \PP\bigg(|x_j|> \frac{r}{2} \text{ for some } j \in [p+1,d]\bigg) 
        = & O\bigg( \sum_{j=p+1}^d \exp\bigg( -\frac{\lambda_p^2}{2\lambda_j^2} R^2 \bigg) \bigg) \\
        \leq & O\bigg( \sum_{j=p+1}^d \exp\bigg( -\frac{R^2}{2} \bigg) \bigg)\\
        = & O\bigg( \exp\bigg( -\frac{R^2}{2}+\log(d-p) \bigg) \bigg),
    \end{align*}
    where the inequality comes from $\lambda_j^2 \leq \lambda_p^2$ for $j=p+1,\ldots,d$. Therefore, we have
    %Now we plug in the exponential decay rate of eigenvalue $\lambda_i^2$'s in Assumption \ref{assump:eigen-decay} and thus arrive at
    %\begin{align*}
    %    \PP\bigg(|x_i|> \frac{r}{2} \text{ for some } i \in [p+1,d]\bigg) 
    %    = & O\bigg( \sum_{i=p+1}^d \exp\bigg( -\frac{R^2}{2} \exp(2b(i-p)) \bigg) \bigg) \\
    %\end{align*}
    \begin{equation*}
        \PP(X \notin S(R,r;p) ) = O \Bigg( \bigg[ \exp \bigg(-\frac{2R^2}{3p} + \log\bigg( \frac{3R^2}{p} \bigg) \bigg) \bigg]^{\frac{p}{2}} \Bigg) + O\bigg( \exp\bigg( -\frac{R^2}{2}+\log(d-p) \bigg) \bigg).
    \end{equation*}  
    
    %Let $R^2=r^2\exp\{b'd\}$. Combine \eqref{eq:prob1} and \eqref{eq:prob2} and we will obtain
    %\begin{align*}
         %\PP(X \notin E_p(R,r) ) &=  O\bigg( \bigg(\frac{R^2}{2} \bigg)^{\lfloor \frac{p}{2} \rfloor}\exp\bigg\{-\frac{R^2}{2} \bigg\} \bigg) + O\bigg(  \exp\bigg\{ -\frac{r^2}{a'} \exp\{b' d\} \bigg\} \bigg)\\
         %&= O\bigg(  \exp\bigg\{ -c r^2 \exp\{b' d\} \bigg\} \bigg).
    %\end{align*}
    Next, we compute the covering number of $S(R,r;p)$ using hypercubes with side length $r>0$, which is denoted as $N_r(S(R,r;p))$. Notice that the first-$p$-dimensional hyper-ellipsoid of $S(R,r;p)$ is contained in a $p$-dimensional hyper-rectangle with side length $2\lambda_i R$ for $i=1,\ldots,p$, while only one hypercube is required to cover the $j$-th dimension for $j=p+1,\ldots,d$. With this observation,  we derive the upper bound for $N_r(S(R,r;p))$:
    \begin{align*}
        N_r(S(R,r;p)) \leq \prod_{i=1}^p \bigg( \frac{2\lambda_i R}{r}\bigg) 
        = \prod_{i=1}^p \bigg( \frac{\lambda_i}{\lambda_p}\bigg),
    \end{align*} 
    where the last equality results from our choice of $p$.
    %Now we substitute $R=r\exp\{b'd/2\}$ and utilize the decay rate of eigenvalues. Then we have
    %\begin{align*}
        %N_r(E_p(R,r)) \leq \frac{\pi^{p/2}}{\Gamma(p/2+1)}\exp\{cpd\}   \cdot a^{p/2} \exp\{-bp(p+1)/4\} 
        %= O(\exp\{ cpd\}) .
    %\end{align*}
    %If we further set $r= \exp\{-cd\}$, we will get
    %\begin{equation*}
    %    \PP(X \notin E_p(R,r) ) = O(\exp\{-r^\alpha\}) \text{  and  } N_r(E_p(R,r))= O(r^{-p}).
    %\end{equation*}
    %where $R=\exp\{-(c-b'/2)d\}$ and $c>b'-2$.
    %Therefore, we obtain $ N_r(E_p(R)) = \tilde{O} (r^{-p_0}) $.
\end{proof}


\section{Proof of Theorem \ref{thm:generalization-exp}}\label{sec:dim-gaussian}

\subsection{Generalization error under exponential eigenvalue decay}\label{sec:dim-exp}
    
    Combining the criteria $\lambda_p^{-1}=2R/r$ and the exponential eigenvalue decay in Assumption \ref{assump:eigen-decay-exp}, we have
    \begin{equation*}
        \frac{1}{\mu} \exp(\theta p ) = \frac{2R}{r}.
    \end{equation*}
    Moreover, by Lemma \ref{lemma:gaussian}, we can compute the covering number of $S(R,r,p)$:
    \begin{align*}
        N_r(S(R,r;p)) &\leq \prod_{i=1}^p \bigg( \frac{\lambda_i}{\lambda_p}\bigg) =  \prod_{i=1}^p \exp(\theta(p-i)) \leq \exp\biggl( \frac{\theta p^2}{2} \biggr) =  \biggl( \frac{2\mu R}{r} \biggr)^{p/2},
    \end{align*}
    which indicates that effective Minkowski dimension of $P_{\rm data}$ is at most $p/2$. Let $r=n^{-(1-\eta)/(2\beta+\sqrt{\log n /\theta})}$ and $R=\log n$ where $\eta\in(0,1)$ is an arbitrarily small constant. Then we obtain
    \begin{align*}
        \theta p = \log\biggl( \frac{2\mu R}{r} \biggr)=  \frac{1}{\theta} \log( 2\mu)+ \log( \log n )+ \frac{(1-\eta)\log n}{2\beta +\sqrt{\log n /\theta}}  \leq \frac{2(1-\eta)\log n}{2\beta +\sqrt{\log n /\theta}}.
    \end{align*}
    Thereby, we can compute the probability ourside $S(R,r;p)$:
    \begin{align*}
        \PP(X \notin S(R,r;p) ) =& O \Biggl( \biggl[ \exp \bigg(-\frac{2R^2}{3p} + \log\bigg( \frac{3R^2}{p} \bigg) \bigg) \biggr]^{\frac{p}{2}}  + \exp\biggl( -\frac{R^2}{2}+\log(d-p)  \biggr) \Biggr) = O(n^{-\log n/3}).
    \end{align*}
    
    Apply Theorem \ref{thm:generalization} with our choice of $R$, $r$ and $p$. Accordingly, the tuple $(L,B,K)$ is set as
    \begin{equation*}
        L= C_1,~~B=O\biggl(n^{\frac{\beta s}{2\beta+\sqrt{\log n/\theta}}} (\log n)^{\beta s} \biggr), ~~ \text{and}~K=O\biggl(n^{\frac{\sqrt{\log n/\theta}}{2\beta + \sqrt{\log n/\theta} } } \biggr).
    \end{equation*}
    Then we can get
    \begin{align*}
        \EE \norm{ \hat{f} - f^*}^2_{L^2(P_{\rm data})} &= O\biggl(  \PP(X \notin S(R,r;p) ) + \sigma r^{2\beta} + \frac{\sigma^2}{n} \cdot \biggl(\frac{R_S}{r} \biggr)^{p/2} \log\biggl( \frac{(R_S/r)^{p/2}}{ r^{4\beta}-4\tau}\biggr) \biggr)\\
        &= O\biggl(  \sigma n^{- \frac{2\beta(1-\eta)}{2\beta+\sqrt{\log n/\theta}}} + \frac{\sigma ^2}{n} \cdot n^ \frac{(1-\eta)^2\log n /\theta}{ (2\beta +\sqrt{\log n /\theta})^2} \cdot (\log n)^\frac{(1-\eta)\log n /\theta}{2\beta +\sqrt{\log n /\theta}} \cdot \frac{\log n/\theta}{2\beta +\sqrt{\log n /\theta}}\cdot \log n \biggr)\\
        &= O\biggl( \sigma  n^{- \frac{2\beta(1-\eta)}{2\beta+\sqrt{\log n/\theta}}} + \sigma^2  n^{-1 + \frac{(1-\eta)^2\log n /\theta}{ (2\beta +\sqrt{\log n /\theta})^2} + \frac{(1-\eta)\log(\log n)/\theta}{2\beta +\sqrt{\log n /\theta}} }   \cdot (\log n)^{3/2} \biggr).
    \end{align*}
    The last equality utilizes the fact that $(\log n)^{\log n}=n^{\log(\log n)}$. Furthermore, notice that for sufficiently large $n$ satisfying $\log(\log n) / \sqrt{\theta \log n} \leq \eta$, we have
    \begin{align*}
        \frac{(1-\eta)^2\log n /\theta}{ (2\beta +\sqrt{\log n /\theta})^2} + \frac{(1-\eta)\log(\log n)/\theta}{2\beta +\sqrt{\log n /\theta}} &\leq  \frac{(1-\eta)^2\log n /\theta}{ (2\beta +\sqrt{\log n /\theta})^2} + \frac{(1-\eta)\eta \sqrt{\log n/\theta}}{2\beta +\sqrt{\log n /\theta}} \\
        &\leq    \frac{(1-\eta)^2 \sqrt{\log n /\theta} }{ 2\beta +\sqrt{\log n /\theta}} + \frac{(1-\eta)\eta \sqrt{\log n/\theta}}{2\beta +\sqrt{\log n /\theta}}\\
        &\leq \frac{(1-\eta) \sqrt{\log n/\theta}}{2\beta +\sqrt{\log n /\theta}}.
    \end{align*}
    %\begin{align*}
    %    -1 + \frac{\log n /\theta}{ (2\beta +\sqrt{\log n /\theta})^2} &= - \frac{4\beta^2 + 4\beta\sqrt{\log n /\theta} }{ (2\beta +\sqrt{\log n /\theta})^2} \\
    %    & = -\frac{2\beta}{ 2\beta +\sqrt{\log n /\theta}} \cdot \frac{2\beta + 2 \sqrt{\log n /\theta} }{ 2\beta +\sqrt{\log n /\theta}} \\
    %    &\leq -\frac{2\beta}{ 2\beta +\sqrt{\log n /\theta}}. 
    %\end{align*}
    %Also, we observe that 
    %\begin{align*}
    %     (\log n)^\frac{\log n /\theta}{2\beta +\sqrt{\log n /\theta}} &= \exp\biggl( \frac{\log(\log n)\cdot\log n /\theta}{2\beta +\sqrt{\log n /\theta}} \biggr)\\
    %     &\leq \exp\Bigl( \log(\log n)\cdot\sqrt{\log n /\theta} \Bigr)\\
    %     &\leq  \exp(\eta \log n) = n^\eta
    %\end{align*}
     Therefore, we use the above observation to derive the following upper bound for generalization error:
    \begin{align*}
         \EE \norm{ \hat{f} - f^*}^2_{L^2(P_{\rm data})} = O\Bigl( \sigma^2 n^{- \frac{2\beta(1-\eta)}{2\beta+\sqrt{\log n/\theta}}}(\log n)^{3/2} \Bigr).
    \end{align*}


\subsection{Generalization error under polynomial eigenvalue decay}\label{sec:dim-poly}
    
     Similarly to last section, we firstly combine the criteria $\lambda_p^{-1}=2R/r$ and the polynomial eigenvalue decay in Assumption \ref{assump:eigen-decay-poly},     \begin{equation*}
         p = \biggl( \frac{2\rho R}{r} \biggr)^{1/\omega}.
    \end{equation*}
    Moreover, by Lemma \ref{lemma:gaussian}, we can compute the covering number of $S(R,r,p)$:
    \begin{align*}
        N_r(S(R,r;p)) &\leq \prod_{i=1}^p \bigg( \frac{\lambda_i}{\lambda_p}\bigg) =  \prod_{i=1}^p \biggl(\frac{p}{i}\biggr)^\omega = \biggl(\frac{p^p}{p!}\biggr)^\omega \leq \biggl(\frac{p^p}{p^{p/2}}\biggr)^\omega =  p^{\omega p/2} =  \biggl(\frac{2\rho R}{r}\biggr)^{p/2} ,
    \end{align*}
    which indicates that effective Minkowski dimension of $P_{\rm data}$ is at most $p/2$. Let $r=n^{-1/(2\beta+ n^\kappa)}$ and $R=n^{1/(2\omega\beta+ \omega n^\kappa)}$ with $\kappa = (1+1/\omega)/\omega$. Then we obtain
    \begin{align*}
        p = \biggl( \frac{2\rho R}{r} \biggr)^{1/\omega} =  n^{\frac{(1+1/\omega)/\omega}{2\beta+ n^\kappa}} = n^{\frac{\kappa}{2\beta+ n^\kappa}}.
    \end{align*}
    Thereby, we can compute the probability outside $S(R,r;p)$:
    \begin{align*}
        \PP(X \notin S(R,r;p) ) =& O \Biggl( \biggl[ \exp \bigg(-\frac{2R^2}{3p} + \log\bigg( \frac{3R^2}{p} \bigg) \bigg) \biggr]^{\frac{p}{2}}  + \exp\biggl( -\frac{R^2}{2}+\log(d-p)  \biggr) \Biggr) = O\Bigl(\exp\Bigl(-n^{\frac{2}{(2\omega\beta+ \omega n^\kappa)}}/3\Bigr) \Bigr).
    \end{align*}
    
    Apply Theorem \ref{thm:generalization} with our choice of $R$, $r$ and $p$. Accordingly, the tuple $(L,B,K)$ is set as
    \begin{equation*}
        L= C_1,~~B=O\Bigl(n^{\frac{(1+1/\omega)\beta s}{2\beta+ n^\kappa } } \Bigr), ~~ \text{and}~K=O\Bigl(n^{\frac{(1+1/\omega) n^{\kappa/(2\beta+n^\kappa)}}{4\beta + 2n^\kappa}} \Bigr).
    \end{equation*}
    Then we have
    \begin{equation}\label{eq:poly-bound}
    \begin{split}
        \EE \norm{ \hat{f} - f^*}^2_{L^2(P_{\rm data})} &= O\biggl(  \PP(X \notin S(R,r;p) ) + \sigma  r^{2\beta} + \frac{\sigma^2}{n} \cdot \biggl(\frac{R_S}{r} \biggr)^{p/2} \log\biggl( \frac{(R_S/r)^{p/2}}{ r^{4\beta}-4\tau}\biggr) \biggr)\\
        &= O\biggl(\sigma n^{- \frac{2\beta}{2\beta+ n^\kappa}} + \frac{\sigma ^2}{n} \cdot n^{ \frac{1+1/\omega}{2\beta+ n^\kappa}\cdot\frac{1}{2}{n^{\frac{\kappa}{2\beta+ n^\kappa}}}} \cdot n^{\frac{\kappa}{2\beta+ n^\kappa}} \log n\biggr)\\
        &= O\Bigl( \sigma  n^{- \frac{2\beta}{2\beta+ n^\kappa}} + \sigma^2 n^{-1+ \frac{1+1/\omega}{2\beta+ n^\kappa}\cdot\frac{1}{2}{n^{\frac{\kappa}{2\beta+ n^\kappa}}} + \frac{\kappa}{2\beta+ n^\kappa}} \log n\Bigr).
    \end{split}
    \end{equation}
    Notice that
    \begin{align*}
        \frac{1+1/\omega}{2\beta+ n^\kappa}\cdot\frac{1}{2}{n^{\frac{\kappa}{2\beta+ n^\kappa}}} + \frac{\kappa}{2\beta+ n^\kappa} &= \frac{1+1/\omega}{2\beta+ n^\kappa} \biggl( \frac{1}{2}{n^{\frac{\kappa}{2\beta+ n^\kappa}}} + \frac{1}{\omega}   \biggr) \\
        & \leq \frac{2}{2\beta+ n^\kappa} \biggl( \frac{1}{2}{n^{\frac{\kappa}{2\beta+ n^\kappa}}} +1 \biggr) \\
        &\leq \frac{n^\kappa}{2\beta+ n^\kappa} , 
    \end{align*}
    where the first inequality is due to $\omega>1$. Therefore, plug the above inequality in \eqref{eq:poly-bound}, we derive the following upper bound for generalization error:
    \begin{align*}
         \EE \norm{ \hat{f} - f^*}^2_{L^2(P_{\rm data})} = O\Bigl( \sigma^2 n^{- \frac{2\beta}{2\beta+ n^\kappa}} \log n\Bigr).
    \end{align*}






\section{Auxiliary Lemmas}

In this section, we investigate the probability tail bound of standard Gaussian variable, which is useful for the proof of Lemma \ref{lemma:gaussian}.  At first, we compute the tail bound for multivariate Gaussian variable.
\begin{lemma}\label{lemma:gaussian-tail-1}
    Suppose $Z=(z_1,\ldots,z_p)\sim N(0,I_p)$ is a standard Gaussian variable in $\RR^p$. Then for any $t>0$, we have
    \begin{equation*}
       \PP\big( \|Z\| > t \big) \leq \bigg( \frac{2t^2 +p}{p} \bigg)^{\frac{p}{2}} \exp \bigg(-\frac{t^4}{2t^2+p} \bigg).
    \end{equation*}
\end{lemma}
\begin{proof}
    By the Markov's inequality, for any $\mu\in(0,1/2)$, we have
    \begin{align*}
        \PP\big( \|Z\| > t \big) =& \PP \big(\exp(\mu \|Z\|^2) > \exp(\mu t^2) \big) \\
        \leq & \frac{\EE \exp(\mu \|Z\|^2) }{\exp(\mu t^2)} \\
        =&  \frac{ \prod_{i=1}^p \EE \exp(\mu z_i^2) }{\exp(\mu t^2)} ,
    \end{align*}
    where the last equality comes from the independence of $z_i$'s.
To bound $\EE \exp(\mu z_i^2)$, we first examine the moment generating function of $z_i$: for any $t\in\RR$,
\begin{align*}
    \EE \exp(t z_i) = \int_\RR \exp(tw)\phi(w) \ud w = \exp(t^2/2),
\end{align*}
where $\phi(w)=(2\pi)^{-p/2} \exp(-w^2/2)$ denotes the probability density function of stardard Gaussian. Then multiply $\exp(-t^2/(2\mu))$ on both sides,
\begin{align*}
    \int_\RR \exp\bigg(tw-\frac{t^2}{2\mu} \bigg)\phi(w) \ud w = \exp\bigg(\frac{t^2(\mu-1)}{2\mu} \bigg).
\end{align*}
By integrating both sides with respect to $t$, we have
\begin{equation*}
    \sqrt{2\pi\mu}  \int_\RR \exp\bigg( \frac{\mu w^2}{2} \bigg)\phi(w) \ud w = \sqrt{\frac{2\pi\mu}{1-\mu} },
\end{equation*}
which indicates
\begin{align*}
    \EE \exp(\mu z_i^2) = \EE \exp\bigg( \frac{2\mu z_i^2}{2} \bigg) = \sqrt{\frac{1}{1-2\mu}}.
\end{align*}
Therefore, for any $\mu \in (0,1/2)$, we have
\begin{align*}
    \PP\big( \|Z\| > t \big) \leq (1-2\mu)^{-\frac{p}{2}}\exp(-\mu t^2).
\end{align*}
Let $\mu = t^2/(2t^2+p)$ and thereby we can conclude the proof of the lemma.

\end{proof}


For standard Gaussian in $\RR$, we derive a tighter upper bound in the following lemma.
\begin{lemma}\label{lemma:gaussian-tail}
    Suppose $z\sim N(0,1)$ is a standard Gaussian variable in $\RR$. Then for any $t>0$, we have
    \begin{equation*}
       \PP\big( \norm{z} > t \big) \leq  \exp\bigg(-\frac{1}{2}t^2\bigg).
    \end{equation*}
\end{lemma}
\begin{proof}
    Firstly, for any $t>0$, compute the probability that $z>t$:
    \begin{align*}
        \PP( z > t ) =& \int_t^\infty \frac{1}{\sqrt{2\pi}}\exp\bigg( -\frac{1}{2}z^2 \bigg) \ud z \\
        = & \int_0^\infty \frac{1}{\sqrt{2\pi}}\exp\bigg( -\frac{1}{2}(u+t)^2 \bigg) \ud u \\
        =& \exp\bigg(-\frac{1}{2}t^2 \bigg)\int_0^\infty \exp(-tu)\cdot\frac{1}{\sqrt{2\pi}}\exp\bigg( -\frac{1}{2}u^2 \bigg) \ud u \\
        \leq & \exp\bigg(-\frac{1}{2}t^2 \bigg)\int_0^\infty \frac{1}{\sqrt{2\pi}}\exp\bigg( -\frac{1}{2}u^2 \bigg) \ud u \\
        = & \frac{1}{2}\exp\bigg(-\frac{1}{2}t^2 \bigg).
    \end{align*}
    Then notice that
\begin{align*}
    \PP\big( \norm{z} > t \big) = \PP( z > t ) + \PP( z < -t ) = 2\PP( z > t ).
\end{align*}
Thereby, we can conclude the proof.
\end{proof}